\documentclass[10pt,a4paper]{amsart}
\usepackage[margin=19mm]{geometry}
\usepackage{amsmath,amssymb,mathtools}
\usepackage[scr=rsfs,cal=euler]{mathalfa}
\usepackage{xfrac}
\usepackage[unicode,hidelinks,bookmarks=false]{hyperref}
\newcommand{\ie}{i.e.\ }
\newcommand{\cf}{cf.\ }
\newcommand{\resp}{respectively\ }
\newcommand{\Subsection}{\subsection}
\providecommand{\nobreakdash}{\nobreak\hbox{-}}
\setlength{\emergencystretch}{2em}
\multlinegap0pt
\usepackage{bm}
\usepackage{bbm}
\usepackage{dsfont}
\usepackage{xspace}
\usepackage{cancel}
\usepackage{mathtools}
\usepackage{amsthm}
\makeatletter
\@ifundefined{theorem}{\newtheorem{theorem}{Theorem}}{}
\@ifundefined{lemma}{\newtheorem{lemma}[theorem]{Lemma}}{}
\makeatother
\newcommand{\adj}{Adj}
\newcommand{\floor}[1]{\lfloor #1 \rfloor}
\newcommand{\inc}{Inc}
\newcommand{\lex}{\mathsf{lex}}
\title[First-Order Logic and Twin-Width for Some Geometric Graphs]{First-Order Logic and Twin-Width for Some Geometric Graphs}
\author{Colin Geniet, Gunwoo Kim, Lucas Meijer}
\date{Source: arXiv:2512.21896v1 (CC BY 4.0)}
\begin{document}
\maketitle

\noindent\emph{Source and licence.} First-Order Logic and Twin-Width for Some Geometric Graphs. Version arXiv:2512.21896v1 (CC BY 4.0), distributed under the Creative Commons Attribution 4.0 International licence, as recorded in the arXiv licence metadata for this exact version and shown on the abstract page https://arxiv.org/abs/2512.21896v1. Full rights notice: licenses/arxiv-2512.21896-CC-BY-4.0.txt.\par\smallskip

\noindent\textit{Editorial scope.} Counted unit: the Lemma whose source label is \texttt{lem:incidence-matrix-grid} in the article (source0.tex lines 384--388, source label \texttt{lem:incidence-matrix-grid}), delivered together with the article's own complete original proof of it (source0.tex lines 389--421, one \texttt{proof} environment of 33 lines). No mathematics is written, repaired or re-derived: statement and proof are the article's own text line for line. Required context retained uncounted: none. Every definition, display and label of those lines is retained unchanged. Omitted from this package and not redistributed: 1--383, 422--1611. Nothing inside a retained excerpt is omitted or altered: every retained body is delivered exactly as the article's lines are written, so no omission receipt of any kind accompanies this package. Citations: the retained excerpts carry no citation command at all, so no bibliography entry of the article is transcribed and no reference is redistributed. Reference adaptation: none was needed and none was made; every reference inside the delivered unit resolves inside it, so no unresolved reference or citation is produced. Printed slips kept literal: no printed word, symbol, hypothesis or number of the counted statement or of its proof was altered. Layout and class: the article's own class is \texttt{lipics-v2021} with packages including mathtools, todonotes, comment, tikz, circuitikz, apxproof, xcolor, xspace, scalerel; the wrapper is the lane's pinned \texttt{amsart} editorial wrapper at 10pt on A4, which restates the theorem-like environments the retained text needs and adds the article's own macro definitions that the retained text needs (\texttt{\textbackslash adj}, \texttt{\textbackslash floor}, \texttt{\textbackslash inc}, \texttt{\textbackslash lex}), with extra wrapper packages (bm, bbm, dsfont, xspace, cancel, mathtools). The delivered numbering is the unit's own. Assets: no retained span contains a figure environment, a graphics-inclusion command, a PDF-page inclusion, a table or a TikZ picture; no image file is copied into this package, the package declares no asset dependency and no image file is redistributed with it. Licence: version arXiv:2512.21896v1 of this article is distributed under the Creative Commons Attribution 4.0 International licence, as recorded in the arXiv licence metadata for this exact version and shown on the abstract page https://arxiv.org/abs/2512.21896v1. A full rights notice is delivered at licenses/arxiv-2512.21896-CC-BY-4.0.txt. Characters: every retained excerpt is pure ASCII, so the delivered engine, XeLaTeX with its default Latin Modern text font, reports no missing character.
\begin{lemma}\label{lem:incidence-matrix-grid}
    Let $E \subseteq V^2$ be a binary relation such that $\adj(E,<_V)$ has no $k$-grid.
    Let~$<_\lex$ be the lexicographic ordering of~$E$, i.e.\ $(x,y) <_\lex (x',y')$ if either~$x <_V x'$, or $x = x'$ and $y <_V y'$.
    Then $\inc(E,<_V,<_\lex)$ has no $4k$-grid.
\end{lemma}

\begin{proof}
    For a directed edge $(x,y) \in E$, we write $s(e) = x$ and $t(e) = y$ for its source and target.
    Assume that the $Inc(E,<_V,<_\lex)$ contains an $\ell$-grid $(p_{i,j})_{i,j \in [\ell]}$.
    Each~$p_{i,j}$ is a~`1' in the matrix, corresponding to a vertex~$v_{i,j}$ and an incident edge~$e_{i,j}$. The definition of $\ell$-grid gives the following inequalities for any $i,i',j,j'$:
    \begin{align}
        \label{eq:ord-vtx} i<i' \quad & \text{implies} \quad v_{i,j} <_V v_{i',j'}, \text{ and}  \\
        \label{eq:ord-edge} j<j' \quad & \text{implies} \quad e_{i,j} <_\lex e_{i',j'}.
    \end{align}
    By definition of the lexicographic ordering, condition~\eqref{eq:ord-edge} also gives 
    \begin{equation}
        \label{eq:ord-source}
        j<j' \quad \text{implies} \quad s(e_{i,j}) \le_V s(e_{i',j'}).
    \end{equation}
    Say that the pair~$(i,j)$ is of \emph{source type} if $v_{i,j} = s(e_{i,j})$, and otherwise of \emph{target type}, meaning $v_{i,j} = t(e_{i,j})$.
    Observe that~$(i+1,j)$ and~$(i,j+1)$ cannot both be of source type, as this would imply both $v_{i,j+1} <_V v_{i+1,j}$ by~\eqref{eq:ord-vtx}, and $v_{i+1,j} \le_V v_{i,j+1}$ by~\eqref{eq:ord-source}.
    By grouping the pairs~$(i,j)$ into $2$-by-$2$ squares $(2i,2j),(2i,2j-1),(2i-1,2j),(2i-1,2j-1)$
    and keeping only one pair of target type in each, we obtain a new $\floor{\ell/2}$-grid in which all pairs~$(i,j)$ are of target type.
    Thus explicitly, we have edges~$f_{i,j}$ for~$i,j < \ell/2$ satisfying for any~$i,i',j,j'$ that
    \begin{align}
        \label{eq:ord-tgt} i<i' \quad & \text{implies} \quad t(f_{i,j}) <_V t(f_{i',j'}), \text{ and}  \\
        \label{eq:ord-edge2} j<j' \quad & \text{implies} \quad f_{i,j} <_\lex f_{i',j'}.
    \end{align}
    We claim that for~$j+1<j'$ and $i>1$, we also have $s(f_{i,j}) <_V s(f_{i',j'})$.
    Indeed, we have
    \[ f_{i,j} <_\lex f_{1,j+1} <_\lex f_{i',j'}, \]
    which translates into the corresponding non-strict inequalities on the respective source vertices.
    But on the other hand,~\eqref{eq:ord-tgt} gives $t(f_{i,j}) >_V t(f_{1,j+1})$.
    Then, for $f_{i,j} <_\lex f_{1,j+1}$ to hold, it must be that $s(f_{i,j}) <_V s(f_{1,j+1})$.
    Thus, when considering edges~$f_{i,j}$ only for~$i>1$ and odd values of~$j$,
    we have that the order of sources~$s(f_{i,j})$ coincides with the ordering of $j$-indices,
    and the ordering of targets~$t(f_{i,j})$ with that of $i$-indices.
    This yields an $\floor{\ell/4}$-grid in~$\adj(E,<_V)$, a contradiction.
\end{proof}

\end{document}
